% 第1章 索洛模型与序列视角下的动态优化

% 导言

\chapter{索洛模型与序列视角下的动态优化} \label{ch-solow-dynamic-opt-sq}

现代宏观经济学通过对经济主体行为的优化来刻画宏观经济的动态运行过程。因此，动态优化是学习量化宏观的基础。
宏观经济理论发展至今，对动态优化问题形成了两种处理方法：一是序列问题视角，通过求解差分方程为经济主体选择一个最优的序列（有限或是无穷），
这也是处理动态问题的基础视角；二是递归视角：利用动态规划方法将动态优化问题的解表示成一个函数——将当期的状态映射到下期，从而降低求解的维度。

本章将先对经典的索洛模型进行回顾，引出动态宏观的基本概念和相关理论；
接着，借助现代宏观经济学的基石模型（Workhouse Model）——新古典增长模型（或称最优增长模型），
站在序列问题的视角下对动态优化理论进行介绍，为下一章递归方法和动态规划学习奠定基础。


% 1.1 =============== 索洛模型 ===============
\section{索洛模型} \label{sec-solow}

现代增长理论始于Ramsey（1928）的经典论文。
Solow（1956）和Swan（1956）在其基础上，建立了Solow-Swan模型来解释长期经济增长，建立了新古典增长理论，奠定了现代宏观经济学的基础。
首先回顾这一经典模型。总量生产函数形式为：
\begin{equation*}
    Y_t = F\left(K_t, L_t\right)
\end{equation*}
其中$Y_t$可以理解为经济的GDP。对于函数$F(K, L)$，有如下假设：
\begin{assumption}  \label{as-F-increase}
    $F(K, L)$关于$k$和$L$严格递增
\end{assumption}
\begin{assumption}  \label{as-F-quasiconcave}
    $F(K, L)$关于$(K, L)$严格拟凹
\end{assumption}
\begin{assumption}  \label{as-F-cosntant}
    $F(K, L)$关于$(K, L)$规模报酬不变
\end{assumption}
\begin{assumption}  \label{as-F-zero}
    $F(0, L)=0$
\end{assumption}
\begin{assumption}  \label{as-F-inada}
    \textbf{Inada条件}：$\lim\limits_{K\rightarrow 0}F_1(K, L)=\infty, \quad \lim\limits_{K\rightarrow \infty}F_1(K, L)=0$
\end{assumption}

在基准模型中，假设人口总数和人均工作时间为常数，因此$L_t$为常数。将人口总数和人均工时正则化为1，生产函数中剩下的变量就是$K_t$，
同时可以写出人均形式：
\begin{equation*}
    F(K_t, 1)=F(k_t, 1)\equiv f(k_t)
\end{equation*}
从而生产函数可以表示为：
\begin{equation}
    y_t = f(k_t)
\end{equation}
假定外生储蓄率为$s$,模型的剩余结构如下：
\begin{itemize}
    \item 投资：$i_t = sy_t$
    \item 资本积累：$k_{t+1}=i_t+(1-\delta)k_t = sf(k_t)+(1-\delta)k_t$
    \item 资源约束：$y_t = i_t + c_t$
\end{itemize}


% 1.2 =============== 最优增长模型 ===============
\section{新古典增长模型} \label{sec-ngm}

Solow（1956）和Swan（1956）的工作标志这新古典增长理论的诞生。
但Solow模型的主要问题在于，模型中储蓄率是外生给定的。
Cass（1965）和Koopmans（1965）在Solow工作的基础上，通过引入微观基础对新古典理论进行了拓展，建立了确定性情形下的最优增长模型：
他们利用家庭偏好来解释为何储蓄是一种内生理性选择的结果，同时将资本积累内生化。
最优增长模型通常被称为RCK模型（Ramsey-Cass-Koopmans Model）。

典型的确定性最优增长模型形式如下：
假设时间是离散且无穷的，由$t=0,1,\dots, \infty$表示，给定第0期资本存量：
$$k_0>0$$
中央计划者（Social Planner）在资源约束:
$$c_t+k_{t+1}=f\left(k_t\right)+(1-\delta)k_t,\quad \forall t \geq 0$$
和资源非负的约束下：
$$c_t \geq 0, \quad k_{t+1} \geq 0, \quad \forall t \geq 0$$
选择无穷期序列：
$$\left\{c_t, k_{t+1}\right\}_{t=0}^{\infty}$$
最大化无穷期效用函数：
$$\sum_{t=0}^{\infty} \beta^t u\left(c_t\right)$$

这一形式通常被称为\textbf{序列问题（Sequence Probem）}。同时，模型还有如下的规范化假设，我们将在下一部分具体理解这些假设的用途。
\begin{assumption} \label{as-f}
    生产函数$f\left(k\right)$连续，二次可微
        且$f\left(0\right)=0,\quad f^{\prime}\left(k\right)>0,\quad f^{\prime\prime}\left(k\right)<0$
\end{assumption}
\begin{assumption} \label{as-f-inada}
    生产函数Inada条件：
        $\lim\limits_{k \rightarrow 0} f^{\prime}(k)=\infty, \quad \lim\limits_{k \rightarrow \infty} f^{\prime}(k)=0$
\end{assumption}
\begin{assumption} \label{as-u}
    效用函数$u$连续，二次可微
        且$u^{\prime}\left(c\right)>0,\quad u^{\prime\prime}\left(c\right)<0$
\end{assumption}
\begin{assumption} \label{as-u-inada}
    效用函数Inada条件：
    $\lim\limits_{c \rightarrow 0} u^{\prime}(c)=\infty, \quad \lim\limits_{c \rightarrow \infty} U^{\prime}(c)=0$
\end{assumption}


% 1.3 ===============序列问题：基本Lagrangian方法与横截性条件 ===============
\section{序列问题：基本Lagrangian解法与横截性条件} \label{sec-sp}

利用Lagrangian函数是处理优化问题的基本方法。但当问题拓展到上面无穷期形式时，需要一些额外条件来确保问题的可解性。
我们将从有限期情形开始，逐步推广至无限期的情形。


% 1.3.1 ===== 有限期情形
\subsection{有限期情形} \label{subsec-finite-h}

首先考虑上面问题的有限期形式，
即时间$t \in$ $\{0,1, \ldots T\}$。对于最大化问题，记第$t$期资源约束相关的乘子为$\beta^t \lambda_t$，进而可构造如下拉格朗日方程：
$$
\mathcal{L}=\sum_{t=0}^T \beta^t u\left(c_t\right)-\sum_{t=0}^T \beta^t \lambda_t\left[c_t+k_{t+1}-(1-\delta) k_t-f\left(k_t\right)\right]
$$
也即：
$$
\mathcal{L}=\sum_{t=0}^T \beta^t\left\{u\left(c_t\right)-\lambda_t\left[c_t+k_{t+1}-(1-\delta) k_t-f\left(k_t\right)\right]\right\}
$$
不难看出，$\lambda_t$是计划者第$t$期资源约束的影子价格，或者说是边际价值。在生产函数和效用函数满足规范化假设时，只要$k_t>0$，这一问题有内点解。
分别对$c_t$和$k_{t+1}$求一阶条件，可得：
$$u^{\prime}\left(c_t\right)-\lambda_t=0$$
$$-\lambda_t+\beta\lambda_{t+1}\left[1-\delta+f^{\prime}\left(k_{t+1}\right)\right]=0$$
与乘子$\lambda_t$相关的一阶条件即为资源约束：在第$t$期，如果经济额外获得$\epsilon$足够小单位的商品，那么从第$t$期看，
福利增加了$\lambda_t\epsilon$；从第0期看，则福利增加了$\beta^t\lambda_t$单位。
合并两个一阶条件可得：
$$
\frac{u^{\prime}\left(c_t\right)}{\beta u^{\prime}\left(c_{t+1}\right)}=f^{\prime}\left(k_{t+1}\right)+1-\delta
$$
这一条件表明：计划者的最优选择是使得今明两天消费的MRS与MRT相等

对于资本$k$，一个自然的问题是：
终期$T$时如何选择，即$k_{T+1}$的最优值?由$k_{T+1}$的KT条件:
$$\frac{\partial \mathcal{L}}{\partial k_{T+1}} \geq 0, \quad k_{T+1} \geq 0$$
且满足互补松弛条件。这等价于：
$$
\beta^T \lambda_T \geq 0, \quad k_{T+1} \geq 0, \quad \beta^T \lambda_T k_{T+1}=0
$$
这意味着在终期第$T$期选择的$k_{T+1}=0$或$k_{T+1}$的影子价值为0，因为不存在$T+1$期，从而在第$T$期选择$k_{T+1}$没有意义。
又由于$\lambda_T=u^{\prime}\left(c_T\right)$，而模型\textbf{假设}$u^{\prime} > 0$，因此，终期T存在角点解：
$$k_{T+1}=0$$
或者说，如果终期选择$k_{T+1}>0$，则计划者会余下未使用资源，在规范化假设下这不会是最优策略（因为$u^{\prime}>0$）。

% 1.3.2 ===== 无限期情形
\subsection{无限期情形} \label{subsec-infinit-h}

现在，\textbf{由有限期转向无限期问题}。回忆在有限期问题中，终期时$k_{T+1}$的最优条件满足：
$$\beta^T \lambda_T k_{T+1}=0$$
从直觉出发，令$T \rightarrow \infty$，那么可以得到无限期情形下的\textbf{横截性条件}（Transversality Condition）：
\begin{equation}
    \lim _{T \rightarrow \infty} \beta^T \lambda_T k_{T+1}=0
\end{equation}
\textbf{这一条件排除了无限期情形下资本无限增长的可能性}。计划者的最优配置$\left\{c_t, k_{t+1}\right\}_{t=0}^{\infty}$是符合
如下欧拉方程（optimality）、资源约束（feasibility）和边界条件的一条\textbf{唯一路径}（unique path）：
\begin{equation}    \label{eq-euler-sp}
    \frac{u^{\prime}\left(c_t\right)}{\beta u^{\prime}\left(c_{t+1}\right)}=f^{\prime}\left(k_{t+1}\right)+1-\delta
\end{equation}
\begin{equation}    \label{eq-res-const-sp}
    k_{t+1}=f\left(k_t\right)+(1-\delta) k_t-c_t
\end{equation}
以及两个边界条件
$$
k_0>0, \quad \lim _{t \rightarrow \infty} \beta^t u^{\prime}\left(c_t\right) k_{t+1}=0
$$
式\eqref{eq-euler-sp}被称为\textbf{欧拉方程（Euler Equation）}。
事实上，符合前两条条件的$\left\{c_t, k_{t+1}\right\}_{t=0}^{\infty}$路径可能有多条，但在两个边界条件约束下，这条路径是唯一确定的。后续将会证明。